% Lösungen Potenzen, Zehnerpotenzen und Quadratwurzeln (zusammenbau v0.4, Heft)
\einheitenkopf{Potenzen, Zehnerpotenzen und Quadratwurzeln}

\erg{45}{a) hoch – eine Potenz aus vier gleichen Faktoren \quad b) $2 \cdot 2 \cdot 2 \cdot 2 \cdot 2 \cdot 2 = 64$ \quad c) $(-4) \cdot (-4) \cdot (-4) = -64$}
\erg{46}{a) $2, 4, 8, 16, 32, 64, 128, 256, 512$: neun Faktoren, also $x = 9$ \quad b) $2, 4, 8, 16, 32, 64, 128$: sieben Faktoren, also $x = 7$ \quad c) $4, 16, 64$: drei Faktoren, also $x = 3$ \quad d) $4, 16, 64, 256, 1\,024$: fünf Faktoren, also $x = 5$ \quad e) $9^3 = 729$, $3^5 = 243$, $2^9 = 512$ – die größte ist $9^3$ \quad f) $6^4 = 1\,296$, $4^5 = 1\,024$, $3^6 = 729$ – die größte ist $6^4$ \quad g) $5^{-3} = \frac{1}{125} = 0{,}008 < 0{,}01$ \quad h) $2^{-4} = \frac{1}{16} = 0{,}0625 > 0{,}05$}
\erg{47}{a) größer als eins – die Hochzahl ist positiv, das Komma wandert nach rechts \quad b) $10\,000$ \quad c) $0{,}0036$}
\erg{48}{a) $630\,000$ – das Komma wandert fünf Stellen nach rechts \quad b) $2\,070\,000$ – das Komma wandert sechs Stellen nach rechts \quad c) $9\,100\,000\,000$ – das Komma wandert neun Stellen nach rechts \quad d) $86\,000\,000\,000$ \quad e) $8\,100\,000\,000$ \quad f) $71\,000$}
\erg{49}{a) $5$ – das Komma wandert fünf Stellen \quad b) $5$ – das Komma wandert fünf Stellen \quad c) $0{,}00034$ – das Komma wandert vier Stellen nach links \quad d) $0{,}000068$ – das Komma wandert fünf Stellen nach links \quad e) $51\,200\,000\,000\,000$ Byte – zwei Nullen angehängt \quad f) $280\,000\,000\,000$-mal – zwei Nullen angehängt \quad g) $-4 \cdot 10^{3}$, denn $-4 \cdot 10^{3} = -4\,000$ liegt auf der Zahlengerade am weitesten links \quad h) $-7 \cdot 10^{5}$, denn $-7 \cdot 10^{5} = -700\,000$ liegt auf der Zahlengerade am weitesten links}
\erg{50}{a) ja, Quadratzahl – elf mal elf \quad b) $7$, Probe: $7 \cdot 7 = 49$ \quad c) $(-8)^2 = 64$, $\sqrt{64} = 8$}
\erg{51}{a) $\sqrt{6{,}25} = 2{,}5$, die Seite ist $2{,}5$ m lang \quad b) $\sqrt{1\,600} = 40$, die Seite ist $40$ m lang \quad c) $\sqrt{441} = 21$, die Seite ist $21$ cm lang}
\erg{52}{a) $\sqrt{(-7)^2} = 7$, denn erst das Quadrat: $(-7)^2 = 49$, und die Wurzel ist nie negativ \quad b) $\sqrt{(-13)^2} = 13$, denn erst das Quadrat: $(-13)^2 = 169$, und die Wurzel ist nie negativ \quad c) $\frac{7}{3} > \frac{9}{4}$, denn $\frac{7}{3} \approx 2{,}33$ und $\frac{9}{4} = 2{,}25$; $\sqrt{5} \approx 2{,}24$ ist kleiner \quad d) $\sqrt{3} < \frac{7}{4}$, denn $\sqrt{3} \approx 1{,}73$ und $\frac{7}{4} = 1{,}75$ \quad e) $-0{,}26 < -\frac{1}{4} < 1{,}7 < \sqrt{3}$, denn $-\frac{1}{4} = -0{,}25$ und $\sqrt{3} \approx 1{,}73$ \quad f) $-0{,}62 < -\frac{3}{5} < 3{,}3 < \sqrt{11}$, denn $-\frac{3}{5} = -0{,}6$ und $\sqrt{11} \approx 3{,}32$ \quad g) $\sqrt{7{,}5^2 - 2{,}1^2} = \sqrt{56{,}25 - 4{,}41} = \sqrt{51{,}84} = 7{,}2$, die Leiter reicht $7{,}2$ m hoch \quad h) $\sqrt{4^2 - 3{,}6^2} = \sqrt{16 - 12{,}96} = \sqrt{3{,}04} \approx 1{,}74$, der Höhenunterschied ist etwa $1{,}74$ m \quad i) $r = \sqrt{750 : (\pi \cdot 9)} = \sqrt{26{,}53\ldots} \approx 5{,}15$, der Radius ist etwa $5{,}15$ cm \quad j) $330$ ml $= 330\,\text{cm}^3$; $r = \sqrt{330 : (\pi \cdot 8)} = \sqrt{13{,}13\ldots} \approx 3{,}62$, der Radius ist etwa $3{,}62$ cm \quad k) $x = 4 \pm \sqrt{16 - 13} = 4 \pm \sqrt{3}$, also $x_1 \approx 5{,}73$ und $x_2 \approx 2{,}27$ \quad l) $x = -2 \pm \sqrt{4 + 1} = -2 \pm \sqrt{5}$, also $x_1 \approx 0{,}24$ und $x_2 \approx -4{,}24$}
